\documentclass[11pt]{article}
\usepackage[letterpaper,margin=0.8in]{geometry}
\usepackage[T1]{fontenc}
\usepackage[utf8]{inputenc}
\usepackage{lmodern}
\usepackage{xcolor}
\usepackage{hyperref}
\usepackage{enumitem}
\usepackage{parskip}
\usepackage{array}
\usepackage{tabularx}
\usepackage{booktabs}
\usepackage{fancyhdr}
\usepackage{lastpage}

\hypersetup{
  hidelinks,
  pdfauthor={OpenAI},
  pdftitle={Harness SaaS Continuous Delivery Discovery Handout},
  pdfsubject={Platform engineering interview guide},
  pdfcreator={ChatGPT}
}

\setlength{\parindent}{0pt}
\setlength{\emergencystretch}{3em}
\raggedbottom
\setlength{\headheight}{14pt}

\definecolor{HeadingBlue}{HTML}{183A66}
\definecolor{RuleGray}{HTML}{C8D1DC}
\definecolor{AccentGray}{HTML}{4E5B6A}

\pagestyle{fancy}
\fancyhf{}
\lhead{Harness CD Discovery Handout}
\rhead{Platform Engineering Working Session}
\cfoot{\thepage\ of \pageref*{LastPage}}
\renewcommand{\headrulewidth}{0.4pt}
\renewcommand{\footrulewidth}{0pt}

\newcommand{\doccluster}[1]{%
  {\small\textcolor{AccentGray}{\textbf{Official Harness docs:} #1}}\par\vspace{0.4em}
}

\newcommand{\qblock}[5]{%
  \textbf{#1}\par
  \begin{itemize}[leftmargin=1.5em,itemsep=0.2em,topsep=0.3em]
    \item[--] \textbf{Why it matters:} #2
    \item[--] \textbf{Strong answer:} #3
    \item[--] \textbf{Risky answer:} #4
    \item[--] \textbf{Harness doc topic:} #5
  \end{itemize}
  \vspace{0.4em}
}

\begin{document}

{\LARGE\bfseries Harness SaaS Continuous Delivery Discovery Handout}\par
\vspace{0.25em}
{\large Platform-level meeting questions for a working session with a platform team member}\par
\vspace{0.5em}
\textbf{Scope:} Official Harness documentation only. No blogs, forums, Reddit, or inferred internal process details. Where Harness documentation does not prescribe an organizational decision, this handout treats that gap as a discovery target rather than filling it in.\par
\vspace{0.6em}
\hrule height 0.8pt
\vspace{0.7em}

\section*{Pre-read briefing}
\doccluster{\href{https://developer.harness.io/docs/continuous-delivery/overview}{Overview}; \href{https://developer.harness.io/docs/continuous-delivery/x-platform-cd-features/services/services-overview}{Services overview}; \href{https://developer.harness.io/docs/continuous-delivery/x-platform-cd-features/environments/environment-overview}{Environments overview}; \href{https://developer.harness.io/docs/platform/delegates/delegate-concepts/delegate-overview/}{Overview}; \href{https://developer.harness.io/docs/platform/get-started/overview/}{Overview}; \href{https://developer.harness.io/docs/platform/pipelines/input-sets}{Input sets and overlays}; \href{https://developer.harness.io/docs/category/triggers/}{Triggers}; \href{https://developer.harness.io/docs/platform/templates/template}{Templates overview}; \href{https://developer.harness.io/docs/platform/role-based-access-control/rbac-in-harness/}{Role-based access control (RBAC) in Harness}; \href{https://developer.harness.io/docs/platform/governance/policy-as-code/harness-governance-overview}{Harness Policy As Code overview}.}

\begin{itemize}[leftmargin=1.5em,itemsep=0.25em]
  \item Harness CD focuses on deployment: it pulls pre-built artifacts and deploys them to target infrastructure.
  \item A \textbf{Service} is what you deploy and contains artifacts, manifests or specifications, configuration, and service-specific variables.
  \item An \textbf{Environment} is where you deploy and is categorized as \textbf{prod} or \textbf{non-prod}.
  \item Each Environment contains one or more \textbf{Infrastructure Definitions} that represent the physical targets.
  \item In Harness CD, a stage models \textbf{what} you deploy through Services, \textbf{where} you deploy through Environments, and \textbf{how} you deploy through execution steps.
  \item Harness \textbf{Delegates} run in your network or VPC and execute deployment and integration tasks on behalf of Harness Manager.
  \item \textbf{Git Experience} can manage pipelines, templates, and input sets in Git, but it does not sync every platform object.
  \item \textbf{Input sets} and overlays pre-populate runtime values for repeatable executions.
  \item \textbf{Triggers} can start pipelines from Git events, schedules, artifacts, webhooks, and other supported events.
  \item \textbf{Templates}, \textbf{RBAC}, and \textbf{Policy as Code} are the main Harness-native levers for consistency, control, and governance.
\end{itemize}

\section*{1. CD operating model and ownership}
\doccluster{\href{https://developer.harness.io/docs/continuous-delivery/overview}{Overview}; \href{https://developer.harness.io/docs/continuous-delivery/manage-deployments/deployment-concepts}{Deployment concepts and strategies}; \href{https://developer.harness.io/docs/platform/role-based-access-control/rbac-in-harness/}{Role-based access control (RBAC) in Harness}.}

\qblock{1. Where does CI end and CD begin in our operating model?}{Harness documents a real boundary: CD deploys pre-built artifacts, while CI is the separate build-and-test capability. This question exposes the artifact handoff, evidence expectations, and release control point.}{"CI ends at a versioned immutable artifact plus release evidence. CD starts with artifact selection, environment targeting, approvals, and deployment execution."}{"Anything after merge is CD."}{Deployment concepts and strategies}

\qblock{2. Who owns the shared CD platform layer versus application-team pipeline content?}{Harness supports account, organization, and project scopes, but the ownership split is an internal design choice. This question surfaces centralized versus delegated administration.}{"Platform owns delegates, shared templates, guardrails, RBAC patterns, and core approval policy. Teams own service-specific configuration within approved patterns."}{"Each team manages everything in its project however it wants."}{Role-based access control (RBAC) in Harness}

\section*{2. Pipeline architecture and stage design}
\doccluster{\href{https://developer.harness.io/docs/continuous-delivery/manage-deployments/deployment-concepts}{Deployment concepts and strategies}; \href{https://developer.harness.io/docs/continuous-delivery/cd-onboarding/new-user/cd-pipeline-modeling-overview}{4. Deployment Pipeline Modeling Overview}; \href{https://developer.harness.io/docs/continuous-delivery/x-platform-cd-features/advanced/multiserv-multienv/}{Use multiple services and environments in a deployment}.}

\qblock{3. What is the default stage pattern for a production-capable deployment pipeline?}{Harness lets teams compose stages flexibly. A platform team should still define a default architecture to avoid pipeline drift.}{"We standardize stage order, required gates, and failure handling by workload class."}{"Teams design their own stage sequence from scratch."}{Deployment concepts and strategies}

\qblock{4. When do we allow multi-service or multi-environment deployments instead of one service per deployment path?}{Harness supports these patterns but documents behavior and limitations. This question reveals how much coupling the platform is willing to tolerate.}{"Single-service is the default. Bundled deployments are exceptions with explicit justification and ownership."}{"We group many services into one pipeline because it is simpler to click once."}{Use multiple services and environments in a deployment}

\section*{3. Service, environment, and infrastructure modeling}
\doccluster{\href{https://developer.harness.io/docs/continuous-delivery/x-platform-cd-features/services/services-overview}{Services overview}; \href{https://developer.harness.io/docs/continuous-delivery/x-platform-cd-features/environments/environment-overview}{Environments overview}; \href{https://developer.harness.io/docs/continuous-delivery/overview}{Overview}.}

\qblock{5. What is the rule for defining a Harness Service in this organization?}{Harness says a Service represents what you are deploying, but it does not define your internal taxonomy. This question forces a deployable-unit model.}{"A Service maps to one independently deployable workload with a clear artifact and manifest boundary."}{"A Service is whatever a team happens to call a service."}{Services overview}

\qblock{6. How do we model lifecycle environments, and what makes an environment prod versus non-prod?}{Harness environment type is first-class. The answer should not rely only on naming conventions.}{"Environment type is authoritative and tied to approvals, RBAC, and promotion controls."}{"If the name contains prod, we treat it as production."}{Environments overview}

\qblock{7. When do we add a new Infrastructure Definition versus a new Environment?}{Harness environments can contain multiple Infrastructure Definitions. This question exposes scoping discipline and reuse strategy.}{"A new infrastructure definition is used for another physical target within the same logical environment; a new environment represents a different lifecycle or risk boundary."}{"We create new environments whenever there is another namespace or cluster."}{Overview / Environments overview}

\section*{4. Delegate placement, connectivity, and execution boundaries}
\doccluster{\href{https://developer.harness.io/docs/platform/delegates/delegate-concepts/delegate-overview/}{Overview}; \href{https://developer.harness.io/docs/platform/delegates/manage-delegates/select-delegates-with-selectors}{Use delegate selectors}; \href{https://developer.harness.io/docs/platform/delegates/delegate-concepts/delegate-requirements}{Delegate system requirements}; \href{https://developer.harness.io/docs/platform/reference-architectures/delegate-architecture-bestpractices}{Harness Delegate Architecture Best Practices}.}

\qblock{8. Where do delegates run today, and which network or trust boundary is each one expected to cross?}{Delegates are the execution plane. Their placement determines which infrastructure and integrated systems CD can reach.}{"Delegate placement aligns to trust boundaries, network reachability, and operational ownership."}{"We have a shared delegate that can reach most things."}{Overview}

\qblock{9. What blast radius is acceptable for a delegate, and when do we require delegate selectors?}{Harness supports selectors for targeted execution. This question tests whether routing and access boundaries are deliberate or ad hoc.}{"Selectors are used intentionally for access segregation, not as a troubleshooting patch."}{"We add selectors whenever something fails to find the right delegate."}{Use delegate selectors}

\qblock{10. How do we scale delegates for concurrency, maintenance, and failure?}{Harness documents delegate resource planning and concurrent execution considerations. This question surfaces whether capacity planning is proactive.}{"We have a scaling model, upgrade process, and failure replacement plan tied to concurrency expectations."}{"We add another delegate only after queues and failures become visible."}{Delegate system requirements}

\section*{5. Deployment strategy, promotion flow, and rollback design}
\doccluster{\href{https://developer.harness.io/docs/continuous-delivery/manage-deployments/deployment-concepts}{Deployment concepts and strategies}; \href{https://developer.harness.io/docs/platform/approvals/adding-harness-approval-stages}{Using manual Harness approval stages}; \href{https://developer.harness.io/docs/platform/approvals/custom-approvals}{Adding custom approval stages and steps}.}

\qblock{11. Which deployment strategies are approved for which workload types, and who decides the default?}{Harness supports multiple deployment strategies, but the platform team should define the selection criteria.}{"Rolling, canary, blue-green, or equivalent strategies are chosen by workload type, failure domain, and rollback expectations."}{"Teams pick whatever strategy they are most familiar with."}{Deployment concepts and strategies}

\qblock{12. How does promotion from non-prod to prod work in practice?}{This question exposes whether promotion preserves artifact immutability and controlled inputs, or whether production deploys are effectively rebuilt decisions.}{"Promotion reuses the same tested artifact and controlled configuration set, with approvals at defined boundaries."}{"Prod just reruns the pipeline with different values."}{Deployment concepts and strategies / Input sets and overlays}

\qblock{13. What is our rollback contract for technical failure versus bad-but-successful release outcomes?}{Rollback behavior is one of the clearest maturity signals in a CD operating model.}{"Automatic rollback handles technical failure. A defined manual or automated path exists for post-deployment rollback when the deployment succeeded but the release outcome is bad."}{"Rollback depends on who is on call and whether they remember the last good version."}{Deployment concepts and strategies}

\section*{6. Git, input set, and trigger operating model}
\doccluster{\href{https://developer.harness.io/docs/platform/get-started/overview/}{Overview}; \href{https://developer.harness.io/docs/platform/pipelines/input-sets}{Input sets and overlays}; \href{https://developer.harness.io/docs/platform/git-experience/manage-input-sets-in-simplified-git-experience}{Manage Git-backed input sets}; \href{https://developer.harness.io/docs/category/triggers/}{Triggers}; \href{https://developer.harness.io/docs/platform/pipeline-faq/}{Pipeline FAQs}.}

\qblock{14. Is Git or the Harness UI the source of truth for pipelines, templates, and input sets?}{Harness Git Experience has clear scope boundaries. This question exposes whether the operating model is coherent or split-brain.}{"Git is authoritative for pipeline definitions and related Git-backed objects. Platform-only entities that are not Git-backed have explicit owners and change controls."}{"People edit YAML wherever is fastest."}{Overview / Pipeline FAQs}

\qblock{15. How are input sets standardized so operators do not hand-enter critical values at runtime?}{Harness input sets exist to reduce manual variance and make execution repeatable.}{"We publish approved input sets and overlays by scenario, environment, or team pattern, and review them through code or controlled change."}{"Operators usually type the values when they run the pipeline."}{Input sets and overlays / Manage Git-backed input sets}

\qblock{16. Which events are allowed to trigger CD automatically, and what restrictions apply by environment type?}{Harness supports multiple trigger types. This question reveals whether automation boundaries are deliberately controlled.}{"Broad automation is allowed in non-prod. Production triggers are narrow, policy-backed, and tied to approvals and artifact provenance."}{"Any supported trigger can start any deployment if the pipeline author configures it."}{Triggers}

\section*{7. Governance, approvals, RBAC, and policy enforcement}
\doccluster{\href{https://developer.harness.io/docs/platform/approvals/adding-harness-approval-stages}{Using manual Harness approval stages}; \href{https://developer.harness.io/docs/continuous-delivery/x-platform-cd-features/cd-steps/approvals/using-harness-approval-steps-in-cd-stages}{Using manual Harness approval steps in CD stages}; \href{https://developer.harness.io/docs/platform/role-based-access-control/rbac-in-harness/}{Role-based access control (RBAC) in Harness}; \href{https://developer.harness.io/docs/platform/governance/policy-as-code/harness-governance-overview}{Harness Policy As Code overview}; \href{https://developer.harness.io/docs/platform/governance/policy-as-code/harness-governance-quickstart/}{Policy As Code quickstart}; \href{https://developer.harness.io/docs/platform/harness-platform-faqs}{Harness Platform FAQs}.}

\qblock{17. Where are approval gates mandatory in the promotion path, and which approval mechanism do we standardize on?}{Harness supports approval stages and steps, including manual and custom approvals. The question exposes whether approval placement is platform policy or author preference.}{"Approval points are tied to risk transitions and use a standard mechanism with explicit timeout and rejection behavior."}{"Pipelines have approvals if the author decides they need one."}{Using manual Harness approval stages / Using manual Harness approval steps in CD stages}

\qblock{18. How are production versus non-production boundaries enforced in access control?}{Harness states environment access is controlled by environment type rather than pipeline stage. This is a critical governance boundary.}{"Prod access is isolated through resource groups, roles, and environment-type-aware controls."}{"We mostly depend on pipeline names or stage names to tell users what is production."}{Role-based access control (RBAC) in Harness / Harness Platform FAQs}

\qblock{19. Which controls are enforced by RBAC and which are enforced by Policy as Code?}{Harness offers both. Mature platform teams use them for different problems rather than treating either one as a catch-all control.}{"RBAC controls who can act. Policy as Code controls what configurations and execution conditions are allowed on save or run events."}{"We rely on peer review to catch bad YAML and bad permissions."}{Harness Policy As Code overview / Policy As Code quickstart}

\section*{8. Reuse, templates, scaling across teams, and platform guardrails}
\doccluster{\href{https://developer.harness.io/docs/platform/templates/template}{Templates overview}; \href{https://developer.harness.io/docs/platform/templates/templates-best-practices}{Best practices and guidelines for templates}; \href{https://developer.harness.io/docs/platform/templates/reconcile-pipeline-templates/}{Reconcile pipeline template changes}.}

\qblock{20. What must be templated versus what may remain team-specific?}{Templates are the main Harness-native standardization mechanism. This question exposes whether reuse is optional or enforced.}{"Stage patterns, approval logic, deployment skeletons, and guardrail logic are templated. Teams supply only approved runtime or service-specific values."}{"Templates are available, but teams can ignore them."}{Templates overview}

\qblock{21. How do we version and roll out shared templates safely across teams?}{Harness supports template versioning and reconciliation. This question reveals whether shared change is controlled or disruptive.}{"Templates are versioned intentionally, breaking changes are introduced via stable references or planned migrations, and reconciliation is managed."}{"We edit the shared template and downstream pipelines deal with the fallout."}{Best practices and guidelines for templates / Reconcile pipeline template changes}

\section*{9. Troubleshooting, observability, auditability, and failure handling}
\doccluster{\href{https://developer.harness.io/docs/continuous-delivery/x-platform-cd-features/executions/execution-history}{Pipeline Execution History}; \href{https://developer.harness.io/docs/platform/pipelines/executions-and-logs/view-and-compare-pipeline-executions}{View and compare pipeline executions}; \href{https://developer.harness.io/docs/platform/governance/audit-trail/}{View audit trail}; \href{https://developer.harness.io/docs/platform/harness-platform-faqs}{Harness Platform FAQs}.}

\qblock{22. When a deployment fails, what is the first authoritative place operators go for diagnosis?}{Harness gives teams execution history, compiled YAML views, execution comparisons, logs, and audit data. A mature support model defines the primary diagnostic path.}{"Operators start with execution history and compiled YAML, then compare executions and inspect logs or audit trail as needed."}{"Everyone debugs differently depending on who is available."}{Pipeline Execution History / View and compare pipeline executions}

\qblock{23. What do we expect to be auditable for configuration changes and pipeline execution events, and what is outside that audit trail?}{Harness documents audit behavior and limitations. This question reveals whether the team actually understands its evidence surface.}{"We know which changes and execution events are audited, which audit settings are enabled, and which actions are out of scope such as read-only API requests."}{"Auditability is handled somewhere in the platform."}{View audit trail / Harness Platform FAQs}

\section*{Top 7 must-ask questions for a 30-minute meeting}
\begin{enumerate}[leftmargin=1.5em,itemsep=0.25em]
  \item Where does CI end and CD begin in our operating model?
  \item Is Git or the Harness UI the source of truth for pipelines, templates, and input sets?
  \item What is the rule for defining a Harness Service here?
  \item How do we distinguish and enforce prod versus non-prod boundaries?
  \item Where do delegates run, and what is each delegate allowed to reach?
  \item What is the standard promotion path into production, including approvals?
  \item What is the rollback contract for technical failure and bad release outcomes?
\end{enumerate}

\section*{Top 7 follow-up questions for a deeper architecture session}
\begin{enumerate}[leftmargin=1.5em,itemsep=0.25em]
  \item What is the default stage pattern for a production-capable pipeline?
  \item When are multi-service or multi-environment deployments allowed?
  \item When do we add a new Infrastructure Definition versus a new Environment?
  \item Which deployment strategies are approved for which workload types?
  \item How are input sets and overlays standardized and reviewed?
  \item What must be templated versus team-specific?
  \item Which controls are enforced by RBAC versus Policy as Code?
\end{enumerate}

\section*{Seven red flags to watch for}
\begin{itemize}[leftmargin=1.5em,itemsep=0.25em]
  \item Production is identified mostly by naming convention.
  \item Teams edit pipeline definitions in Git and in the UI without a clear source of truth.
  \item A single delegate or delegate pool has broad reach across unrelated trust zones.
  \item Rollback is manual, undocumented, or depends on tribal knowledge.
  \item Approval placement is left entirely to pipeline authors.
  \item Templates are optional guidance rather than an enforcement mechanism.
  \item Nobody can describe the artifact handoff between CI completion and CD execution.
\end{itemize}

\section*{Official Harness documentation referenced}
{\small
\begin{enumerate}[leftmargin=1.5em,itemsep=0.2em]
  \item \href{https://developer.harness.io/docs/continuous-delivery/overview}{Overview}
  \item \href{https://developer.harness.io/docs/continuous-delivery/manage-deployments/deployment-concepts}{Deployment concepts and strategies}
  \item \href{https://developer.harness.io/docs/continuous-delivery/cd-onboarding/new-user/cd-pipeline-modeling-overview}{4. Deployment Pipeline Modeling Overview}
  \item \href{https://developer.harness.io/docs/continuous-delivery/x-platform-cd-features/services/services-overview}{Services overview}
  \item \href{https://developer.harness.io/docs/continuous-delivery/x-platform-cd-features/environments/environment-overview}{Environments overview}
  \item \href{https://developer.harness.io/docs/continuous-delivery/x-platform-cd-features/advanced/multiserv-multienv/}{Use multiple services and environments in a deployment}
  \item \href{https://developer.harness.io/docs/platform/delegates/delegate-concepts/delegate-overview/}{Overview}
  \item \href{https://developer.harness.io/docs/platform/delegates/manage-delegates/select-delegates-with-selectors}{Use delegate selectors}
  \item \href{https://developer.harness.io/docs/platform/delegates/delegate-concepts/delegate-requirements}{Delegate system requirements}
  \item \href{https://developer.harness.io/docs/platform/reference-architectures/delegate-architecture-bestpractices}{Harness Delegate Architecture Best Practices}
  \item \href{https://developer.harness.io/docs/platform/get-started/overview/}{Overview}
  \item \href{https://developer.harness.io/docs/platform/pipelines/input-sets}{Input sets and overlays}
  \item \href{https://developer.harness.io/docs/platform/git-experience/manage-input-sets-in-simplified-git-experience}{Manage Git-backed input sets}
  \item \href{https://developer.harness.io/docs/category/triggers/}{Triggers}
  \item \href{https://developer.harness.io/docs/platform/pipeline-faq/}{Pipeline FAQs}
  \item \href{https://developer.harness.io/docs/platform/approvals/adding-harness-approval-stages}{Using manual Harness approval stages}
  \item \href{https://developer.harness.io/docs/continuous-delivery/x-platform-cd-features/cd-steps/approvals/using-harness-approval-steps-in-cd-stages}{Using manual Harness approval steps in CD stages}
  \item \href{https://developer.harness.io/docs/platform/approvals/custom-approvals}{Adding custom approval stages and steps}
  \item \href{https://developer.harness.io/docs/platform/role-based-access-control/rbac-in-harness/}{Role-based access control (RBAC) in Harness}
  \item \href{https://developer.harness.io/docs/platform/governance/policy-as-code/harness-governance-overview}{Harness Policy As Code overview}
  \item \href{https://developer.harness.io/docs/platform/governance/policy-as-code/harness-governance-quickstart/}{Policy As Code quickstart}
  \item \href{https://developer.harness.io/docs/platform/templates/template}{Templates overview}
  \item \href{https://developer.harness.io/docs/platform/templates/templates-best-practices}{Best practices and guidelines for templates}
  \item \href{https://developer.harness.io/docs/platform/templates/reconcile-pipeline-templates/}{Reconcile pipeline template changes}
  \item \href{https://developer.harness.io/docs/continuous-delivery/x-platform-cd-features/executions/execution-history}{Pipeline Execution History}
  \item \href{https://developer.harness.io/docs/platform/pipelines/executions-and-logs/view-and-compare-pipeline-executions}{View and compare pipeline executions}
  \item \href{https://developer.harness.io/docs/platform/governance/audit-trail/}{View audit trail}
  \item \href{https://developer.harness.io/docs/platform/harness-platform-faqs}{Harness Platform FAQs}
\end{enumerate}
}

\end{document}
